% ============================================================
\chapter{System Design}
\label{ch:system-design}
% ============================================================

\section{Requirement Specification}
\label{sec:requirement-specification}

\subsection{Functional Requirements}
\label{subsec:functional-requirements}

The functional requirements for the prepared system are:

\begin{enumerate}
    \item The system must be able to capture game state in real time from
    emulator RAM.

    \item It must convert raw RAM coordinates into a compact spatial grid
    suitable for NEAT.

    \item It must evolve genomes through mutation, crossover, and speciation.

    \item It must evaluate agent performance using a fitness function based on
    horizontal progress and stagnation penalties.

    \item It must log generation statistics and generate PDF reports
    automatically.
\end{enumerate}


\subsection{Non-Functional Requirements}
\label{subsec:non-functional-requirements}

These are essential for the better performance and stability of the system. The
points below focus on the non-functional requirements of the prepared system.

\subsubsection{Reliability}

The system should recover from interruptions through checkpointing and
persistent CSV logs.

\subsubsection{Maintainability}

The system is broken down into modular sub-components such as the RAM scraper,
spatial quantizer, NEAT engine, logging, and reporting modules.

\subsubsection{Performance}

The object detection equivalent in this project is low-latency state extraction
and action injection, so the pipeline must keep inference overhead minimal.

\subsubsection{Frame Rate and Throughput}

For smooth operation, the system must support efficient frame-by-frame
execution and rapid generation turnover during training.


% ============================================================
\section{Data Flow Diagram}
\label{sec:dfd}

\subsection{DFD Level 0}
\label{subsec:dfd-level-0}

\begin{figure}[H]
\centering
\includegraphics[width=0.95\textwidth]{dfd_level0}
\caption{DFD Level 0}
\label{fig:dfd-level-0}
\end{figure}


\subsection{DFD Level 1}
\label{subsec:dfd-level-1}

\begin{figure}[H]
\centering
\includegraphics[width=0.92\textwidth]{dfd_level1}
\caption{DFD Level 1}
\label{fig:dfd-level-1}
\end{figure}



% ============================================================
\section{Use Case Diagram}
\label{sec:use-case-diagram}

% Tall portrait image: constrained by height so it fits within the text block.
\begin{figure}[H]
\centering
\includegraphics[height=0.78\textheight,width=\textwidth,keepaspectratio]{use_case_diagram}
\caption{Use Case Diagram}
\label{fig:use-case-diagram}
\end{figure}



% ============================================================
\section{Class Diagram}
\label{sec:class-diagram}

% Tall portrait image: constrained by height so it fits within the text block.
\begin{figure}[H]
\centering
\includegraphics[height=0.8\textheight,width=\textwidth,keepaspectratio]{class_diagram}
\caption{Class Diagram}
\label{fig:class-diagram}
\end{figure}



% ============================================================
\section{Sequence Diagram}
\label{sec:sequence-diagram}

\begin{figure}[H]
\centering
\includegraphics[width=\textwidth]{sequence_diagram}
\caption{Sequence Diagram of a Single Frame Evaluation}
\label{fig:sequence-diagram}
\end{figure}



% ============================================================
\section{Component Diagram}
\label{sec:component-diagram}

% Tall portrait image: constrained by height so it fits within the text block.
\begin{figure}[H]
\centering
\includegraphics[height=0.82\textheight,width=\textwidth,keepaspectratio]{component_diagram}
\caption{Component Diagram}
\label{fig:component-diagram}
\end{figure}



% ============================================================
\section{Deployment Diagram}
\label{sec:deployment-diagram}

\begin{figure}[H]
\centering
\includegraphics[width=0.85\textwidth]{deployment_diagram}
\caption{Deployment Diagram}
\label{fig:deployment-diagram}
\end{figure}



% ============================================================
\section{Gantt Chart}
\label{sec:gantt-chart}

The project schedule was tracked using Asana as a Kanban board, organizing
tasks into sprints. Figure~\ref{fig:gantt-chart} illustrates the timeline of
tasks spanning the entire duration of the project.

\begin{figure}[H]
\centering
\includegraphics[width=\textwidth]{gantt_chart}
\caption{Project Gantt Chart}
\label{fig:gantt-chart}
\end{figure}
